\documentclass[10pt]{article}
\usepackage[T1]{fontenc}
\usepackage[utf8]{inputenc}
\usepackage[french]{babel}
\usepackage{amsmath,amssymb,mathrsfs,mathpazo}
\usepackage{geometry}
\geometry{a4paper,margin=2cm}
\usepackage{graphicx,array,tabularx,tabularray,colortbl,multirow,multicol,booktabs,enumitem,pifont,ulem,fancybox,diagbox,cancel}
\usepackage{pstricks,pstricks-add,pst-plot,pst-tree,pst-node,pst-text,pst-func,pst-math,pst-bspline,pst-3dplot,pst-eucl,pst-solides3d}
\let\pscalculate\undefined
\usepackage{pst-bezier}
\usepackage{tikz,pgfplots,tkz-tab}
\usetikzlibrary{arrows,arrows.meta,patterns,calc,positioning,trees,decorations.markings,shapes.geometric,shapes.misc,shapes.arrows,intersections,decorations.pathmorphing,decorations.pathreplacing,backgrounds,shadows,babel}
\pgfplotsset{compat=1.18}
\usepgfplotslibrary{fillbetween}
\usepackage[np]{numprint}
\usepackage{listings,makecell,pas-tableur,twoopt}
\let\barre\undefined
\usepackage[pstricks]{tablvar}
\definecolor{darkgreen}{RGB}{0,100,0}
\providecommand{\red}{\color{red}}
\providecommand{\blue}{\color{blue}}
\providecommand{\green}{\color{green}}
\providecommand{\black}{\color{black}}
\providecommand{\white}{\color{white}}
\providecommand{\gray}{\color{gray}}
\providecommand{\cyan}{\color{cyan}}
\providecommand{\magenta}{\color{magenta}}
\providecommand{\yellow}{\color{yellow}}
\providecommand{\violet}{\color{violet}}
\providecommand{\orange}{\color{orange}}
\providecommand{\euro}{\texteuro}
\providecommand{\EUR}{\texteuro}
\providecommand{\ieme}{\textsuperscript{e}}
\providecommand{\degres}{\ensuremath{^\circ}}
\providecommand{\degree}{\ensuremath{^\circ}}
\providecommand{\vv}[1]{\overrightarrow{#1}}
\providecommand{\Coord}[1]{#1}
\providecommand{\decosix}{}
\providecommand{\starredbullet}{\ensuremath{\bullet}}
\pagestyle{empty}
\setlength{\parindent}{0pt}
\usepackage{adjustbox}\DeclareUnicodeCharacter{2236}{:}\DeclareUnicodeCharacter{2190}{\ensuremath{\leftarrow}}
\providecommand{\R}{\mathbb{R}}
\providecommand{\N}{\mathbb{N}}
\providecommand{\Z}{\mathbb{Z}}
\providecommand{\Q}{\mathbb{Q}}
\providecommand{\D}{\mathbb{D}}
\providecommand{\C}{\mathbb{C}}
\providecommand{\e}{{\mathrm{e}}}
\providecommand{\vect}[1]{\overrightarrow{#1}}
\providecommand{\vectt}[1]{\overrightarrow{\mathrm{#1}}}
\providecommand{\vvt}[1]{\overrightarrow{\mathrm{#1}}}
\providecommand{\np}[2][]{{#2~#1}}
\providecommand{\vv}[1]{{\overrightarrow{#1}}}
\renewcommand{\d}{{\mathrm{d}}}
\providecommand{\degres}{{^\circ}}
\providecommand{\degre}{{^\circ}}
\providecommand{\blue}{\color{blue}}
\providecommand{\red}{\color{red}}
\providecommand{\enskip}{\quad}
\providecommand{\llbracket}{\lbrack\!\lbrack}
\providecommand{\rrbracket}{\rbrack\!\rbrack}
\providecommand{\euro}{{\text{€}}}
\providecommand{\no}{{\text{n°}}}
\providecommand{\interval}[2]{{[#1;#2]}}
\providecommand{\intervFO}[2]{{[#1;#2[}}
\providecommand{\dotfill}{\ldots}

\providecommand{\ds}{\displaystyle}
\providecommand{\textding}[1]{#1}
\providecommand{\Oij}{\(\mathrm{O};\vec{\imath},\vec{\jmath}\)}
\providecommand{\Oijk}{\(\mathrm{O};\vec{\imath},\vec{\jmath},\vec{k}\)}
\providecommand{\Ouv}{\(\mathrm{O};\vec{u},\vec{v}\)}
\providecommand{\intervFF}[2]{\ensuremath{[#1;#2]}}
\renewcommand{\euro}{\ifmmode\text{\texteuro}\else\texteuro\fi}
\DeclareUnicodeCharacter{00B0}{\ifmmode^{\circ}\else\textdegree\fi}
\DeclareUnicodeCharacter{20AC}{\ifmmode\text{\texteuro}\else\texteuro\fi}
% Énoncés : Éducation nationale. Transcription et corrigés : APMEP, Denis Vergès et contributeurs. Découpage et classement : CoopMaths, Sébastien Lozano, Éric Elter et contributeurs. Adaptation de la présentation : Maths Explorer.
% CC BY-SA : https://coopmaths.fr/www/about/
% https://coopmaths.fr/alea/static/bac/2025/tex/bac_2025_06_sujet2_polynesie_3_cor.tex
\begin{document}
\noindent\textit{APMEP / CoopMaths. CC BY-SA. Présentation adaptée par Maths Explorer.}\par\medskip
\begin{enumerate}
	\item On détermine les coordonnées des vecteurs $\vect{\mathrm{AB}}$ et $\vect{\mathrm{AC}}$ : \quad $\vect{\mathrm{AB}}\begin{pmatrix}-2\\1\\5 \end{pmatrix}$
	 et $\vect{\mathrm{AC}}\begin{pmatrix}-1\\-2\\0 \end{pmatrix}$.

	 Ces vecteurs ont des coordonnées qui sont clairement non proportionnelles, donc les vecteurs ne sont pas colinéaires, et donc les trois points A, B et C ne sont pas alignés : ils déterminent bien un plan.

	 \item Comme le repère est orthonormal, en utilisant les coordonnées des vecteurs précédents, qui sont non nuls, on va calculer le produit scalaire des vecteurs $\vect{\mathrm{AB}}$ et $\vect{\mathrm{AC}}$ :

$\vect{\mathrm{AB}}\cdot\vect{\mathrm{AC}} = -2 \times (-1) + 1  \times  (-2) + 5  \times 0 = 2 + -2 + 0 = 0$.

Les vecteurs $\vect{\mathrm{AB}}$ et $\vect{\mathrm{AC}}$ sont non nuls, mais leur produit scalaire l'est, donc les droites (AB) et (AC) sont orthogonales, et donc le triangle ABC est bien rectangle en A.

	\item 
		\begin{enumerate}
		\item Calculons le produit scalaire de $\vect{u}$ avec $\vect{\mathrm{AB}}$ et $\vect{\mathrm{AC}}$, ces deux vecteurs formant une base du plan (ABC).

\begin{itemize}
\item $\vect{u}\cdot \vect{\mathrm{AB}} = 2  \times (-2) + (-1) \times 1 + 1  \times 5 = -4 - 1 + 5 = 0$;
\item $\vect{u}\cdot \vect{\mathrm{AC}} = 2  \times (-1) + (-1) \times (-2) + 1  \times 0 = -2 + 2 + 0 = 0$;
\end{itemize}

 $\vect{u}$ est donc orthogonal à deux vecteurs formant une base de (ABC), il est donc orthogonal à (ABC).

 		Toute droite dirigée par $\vect{u}$, et donc $\Delta$, est donc orthogonale à (ABC).

 		\item $\vect{u}$ est donc un vecteur normal au plan (ABC), d'après la question précédente.

Par propriété, on en déduit qu'une équation de (ABC) est de la forme :

$2x - y + z + d = 0$\quad où $d$ est un réel.

 $\begin{aligned}
 \mathrm{A} \in (\mathrm{ABC})
 &\iff 2x_{\mathrm{A}} - y_{\mathrm{A}} + z_{\mathrm{A}} + d = 0\\
 &\iff 2\times 1 - 3 + 0 + d = 0\\
 &\iff -1 + d = 0\\
 &\iff d = 1
\end{aligned}$

Une équation de (ABC) est donc bien \quad $2x - y + z + 1 = 0$.

 		\item $\Delta$ contient $\mathrm{D}(-2~;~2~;~1)$ et est dirigée par $\vect{u}$, donc, par propriété, une représentation paramétrique de $\Delta$ est :\quad $\begin{cases}
 			x=-2 + 2t\\
 			y= 2 - t\\
 			z= 1 + t\\
 		\end{cases}$\quad avec $t \in \R$
	 \end{enumerate}
 \item Considérons $M_t$, le point de paramètre $t$ sur la droite $\Delta$.

 $\begin{aligned}
		M_t \in (\mathrm{ABC})
		&\iff 2x_{M_t} - y_{M_t} + z_{M_t} + 1 = 0\\
		&\iff 2(-2+2t) - (2 - t) + (1 + t) + 1 = 0\\
		&\iff -4 + 4t - 2 + t + 1 + t + 1 = 0\\
		&\iff 6t - 4 = 0\\
		&\iff t = \dfrac{4}{6} = \dfrac{2}{3}
 \end{aligned}$

$\Delta$ n'a qu'un seul point sur (ABC) (c'est normal, comme elle est orthogonale au plan, elle est sécante au plan), c'est le point de paramètre $\dfrac{2}{3}$ dans la représentation, c'est-à-dire que c'est le point de coordonnées $\left(-2 + 2\times \dfrac{2}{3}~;~2 - \dfrac{2}{3} ; 1 + \dfrac{2}{3}\right)$, soit $\left(-\dfrac{2}{3}~;~\dfrac{4}{3}~;~\dfrac{5}{3}\right)$.

On reconnaît les coordonnées du point H. H est donc l'intersection de $\Delta$ et (ABC), c'est donc l'intersection du plan (ABC) avec la droite passant par D et orthogonale à (ABC) : H est le projeté orthogonal de D sur (ABC).

	\item 
		\begin{enumerate}
		\item On est dans un repère orthonormé :

$\begin{aligned}
	\mathrm{DH} &= \sqrt{\left(-\dfrac{2}{3} - (-2)\right)^2+\left(\dfrac{4}{3} - 2\right)^2+\left(\dfrac{5}{3} - 1\right)^2 } = \sqrt{\left(\dfrac{4}{3}\right)^2+\left(\dfrac{-2}{3}\right)^2+\left(\dfrac{2}{3}\right)^2 } \\
	&= \sqrt{\dfrac{16}{9} + \dfrac{4}{9} + \dfrac{4}{9}} = \sqrt{\dfrac{24}{9}} = \dfrac{\sqrt{24}}{3} = \dfrac{\sqrt{4\times 6}}{3} = \dfrac{2\sqrt{6}}{3}
\end{aligned}$

		On arrive bien à la distance annoncée : \quad $\mathrm{DH} = \dfrac{2\sqrt{6}}{3}$.

		\item Pour le tétraèdre ABCD, on va choisir comme base le triangle ABC, rectangle en A, et la hauteur correspondante, est donc la distance de H au plan (ABC), c'est-à-dire la distance DH.

		Pour calculer $B$, l'aire de la base, on va utiliser comme base de ABC la longueur AB, et comme hauteur correspondante, la longueur AC (car ABC est rectangle en A).

		Donc :\quad $\begin{aligned}[t]
			B &= \dfrac{\mathrm{AB} \times \mathrm{AC}}{2} = \dfrac{\sqrt{(-2)^2 + 1^2 + 5^2}\times \sqrt{(-1)^2 + (-2)^2 + 0^2}}{2} = \dfrac{\sqrt{30}\times \sqrt{5}}{2}= \dfrac{\sqrt{6\times5^2}}{2}\\
			&=\dfrac{5\sqrt{6}}{2}
		\end{aligned}$

		Le volume du tétraèdre est donc :\quad $\begin{aligned}[t]
			V &= \dfrac{1}{3}\times B \times h = \dfrac{1}{3} \times \dfrac{5\sqrt{6}}{2}\times \dfrac{2\sqrt{6}}{3} \\
			&= \dfrac{5\times2\times6}{3\times6} = \dfrac{10}{3}
		\end{aligned}$

		Le volume du tétraèdre est donc de $\dfrac{10}{3} \approx 3,3$.
		\end{enumerate}

	\item Par lecture de la représentation paramétrique de $d$, on peut dire que la droite $d$ est dirigée par : $\vect{u'}\begin{pmatrix} -2 \\ -3 \\ 1\end{pmatrix}$.

	Comme le repère est orthonormé, on a :\quad $\begin{aligned}[t]
		\vect{u} \cdot \vect{u'} &= 2 \times (-2) + (- 1) \times -3 + 1  \times 1\\
		& = -4 + 3 + 1 = 0
	\end{aligned}$

	$\vect{u'}$ est orthogonal à $\vect{u}$ qui est lui-même un vecteur normal à (ABC). On en déduit que $\vect{u'}$ est un vecteur du plan (ABC). Ainsi, la droite $d$ est une droite parallèle au plan (ABC).

	\emph{Remarque :} bien que ça ne soit pas demandé ici, on peut préciser : $d$ passe par le point E de coordonnées $(1 ; 0 ; 1)$. Ces coordonnées ne vérifiant pas l'équation de (ABC), on en déduit que E n'appartient pas au plan, et donc que $d$ est strictement parallèle à (ABC) (elle n'est pas incluse dans le plan).
	\end{enumerate}

\end{document}
